\documentclass[a4,11pt]{aleph-notas}
% Se puede ver la documentación aquí:
% https://github.com/alephsub0/LaTeX_aleph-notas

% -- Paquetes adicionales
\usepackage{enumitem}
\usepackage{url}
\usepackage{array}
\usepackage{booktabs}
\usepackage{longtable}
\usepackage{aleph-comandos}
\hypersetup{
    urlcolor=blue,
    linkcolor=blue,
}


% -- Datos
\institucion{Facultad de Ciencias Exactas, Naturales y Ambientales}
\carrera{Catálogo STEM}
\asignatura{Matemática III}
\tema{Plan tema 11: Integrales dobles}
\autor{Andrés Merino}
\fecha{Periodo 2026-1}

\logouno[0.16\textwidth]{Logos/logoPUCE_04_ac}
\definecolor{colortext}{HTML}{1957d1}
\definecolor{colordef}{HTML}{1957d1}
\fuente{montserrat}


\begin{document}

\encabezado

%%%%%%%%%%%%%%%%%%%%%%%%%%%%%%%%%%%%%%%%
\section*{Resultado de Aprendizaje}
%%%%%%%%%%%%%%%%%%%%%%%%%%%%%%%%%%%%%%%%

%%%%%%%%%%%%%%%%%%%%%%%%%%%%%%%%%%%%%%%%
\subsection*{RdA de la asignatura:}
\begin{itemize}[leftmargin=*]
    \item \textbf{RdA 1:} Identificar conceptos, teoremas y propiedades de funciones vectoriales, derivadas parciales e integral vectorial.
    \item \textbf{RdA 2:} Calcular derivadas parciales e integrales con asistentes computacionales.
    \item \textbf{RdA 3:} Aplicar Cálculo Vectorial para resolver problemas reales.
\end{itemize}

%%%%%%%%%%%%%%%%%%%%%%%%%%%%%%%%%%%%%%%%
\subsection*{RdA del tema:}
\begin{itemize}[leftmargin=*]
    \item Interpretar la integral doble como límite de sumas de Riemann y reconocer sus propiedades fundamentales.
    \item Aplicar el Teorema de Fubini para calcular integrales dobles sobre rectángulos mediante integrales iteradas.
    \item Identificar y describir regiones elementales de tipo I, tipo II y tipo III en $\R^2$.
    \item Calcular integrales dobles sobre regiones elementales y aplicarlas al cálculo de áreas y volúmenes.
\end{itemize}

%%%%%%%%%%%%%%%%%%%%%%%%%%%%%%%%%%%%%%%%
\section*{Introducción}
%%%%%%%%%%%%%%%%%%%%%%%%%%%%%%%%%%%%%%%%

\paragraph{Pregunta inicial:}
Imagine que necesita conocer la cantidad total de lluvia acumulada en una ciudad rectangular después de una tormenta, sabiendo que la intensidad de lluvia varía de un punto a otro. ¿Cómo calcularía esa cantidad total si tiene una fórmula que describe la intensidad en cada punto del mapa?

%%%%%%%%%%%%%%%%%%%%%%%%%%%%%%%%%%%%%%%%
\section*{Desarrollo}
%%%%%%%%%%%%%%%%%%%%%%%%%%%%%%%%%%%%%%%%

%%%%%%%%%%%%%%%%%%%%%%%%%%%%%%%%%%%%%%%%
\subsection*{Actividad 1: Integral doble sobre un rectángulo}
Se introduce la integral doble como límite de sumas de Riemann bidimensionales y se enuncia el Teorema de Fubini para reducirla a integrales iteradas sobre rectángulos.

\paragraph{¿Cómo lo haremos?}
\begin{itemize}[leftmargin=*]
    \item \textbf{Motivación:} retomar la pregunta inicial; discutir intuitivamente la suma de Riemann en dos variables como suma de volúmenes de prismas rectangulares.
    \item \textbf{Clase magistral:} se definen partición, subrectángulos, suma de Riemann e integral doble; se enuncia la integrabilidad de funciones continuas; se presenta el Teorema de Fubini y las propiedades de linealidad y monotonía. Se usa el resumen \href{https://andres-merino.github.io/MatematicaIII-05-N0051/01Resumen/Resumen11.pdf}{Resumen11.pdf}.
    \item \textbf{Implementación computacional:} calcular integrales dobles sobre rectángulos con un asistente computacional; visualizar geométricamente el volumen bajo una superficie.
    \item \textbf{Resolución de ejercicios:} aplicar el Teorema de Fubini a integrales dobles sobre rectángulos usando \href{https://andres-merino.github.io/MatematicaIII-05-N0051/04Ejercicios/Ejercicios11.pdf}{Ejercicios11.pdf}.
\end{itemize}

\paragraph{Verificación de aprendizaje:}
\begin{itemize}[leftmargin=*]
    \item Calcule $\displaystyle\iint_R (2x+3y)\,dA$ donde $R=[0,1]\times[0,2]$. ¿El orden de integración afecta el resultado?
    \item Enuncie el Teorema de Fubini con sus hipótesis. ¿Por qué es importante que $f$ sea acotada?
    \item ¿Qué representa geométricamente $\iint_R f\,dA$ cuando $f(x,y)\geq 0$?
\end{itemize}

%%%%%%%%%%%%%%%%%%%%%%%%%%%%%%%%%%%%%%%%
\subsection*{Actividad 2: Integral doble sobre regiones elementales}
Se extiende la integral doble a regiones no rectangulares mediante las regiones elementales de tipo I y tipo II, y se aplica al cálculo de áreas y volúmenes.

\paragraph{¿Cómo lo haremos?}
\begin{itemize}[leftmargin=*]
    \item \textbf{Motivación:} plantear el problema de integrar sobre una región triangular o circular; discutir la extensión de $f$ por cero fuera de $\Omega$ como puente conceptual desde el rectángulo.
    \item \textbf{Clase magistral:} se definen regiones de tipo I, tipo II y tipo III; se enuncia el teorema de reducción a integrales iteradas para cada tipo; se destaca la importancia de identificar correctamente los límites de integración según el tipo de región. Se usa el resumen \href{https://andres-merino.github.io/MatematicaIII-05-N0051/01Resumen/Resumen11.pdf}{Resumen11.pdf}.
    \item \textbf{Resolución de ejercicios:} calcular integrales dobles sobre regiones elementales y calcular áreas y volúmenes usando \href{https://andres-merino.github.io/MatematicaIII-05-N0051/04Ejercicios/Ejercicios11.pdf}{Ejercicios11.pdf}.
\end{itemize}

\paragraph{Verificación de aprendizaje:}
\begin{itemize}[leftmargin=*]
    \item Describa la región $\Omega$ acotada por $y=x^2$ e $y=x$ como región de tipo I y de tipo II. ¿Cuál orden de integración resulta más conveniente?
    \item ¿Cómo se usa la integral doble para calcular el área de una región plana $\Omega$?
    \item Calcule $\iint_\Omega x\,dA$ donde $\Omega$ es el triángulo con vértices $(0,0)$, $(1,0)$ y $(0,1)$.
\end{itemize}

%%%%%%%%%%%%%%%%%%%%%%%%%%%%%%%%%%%%%%%%
\section*{Cierre}
%%%%%%%%%%%%%%%%%%%%%%%%%%%%%%%%%%%%%%%%

\paragraph{Tarea:}
Resolver del libro \href{https://catalogobiblioteca.puce.edu.ec/cgi-bin/koha/opac-detail.pl?biblionumber=83405&query_desc=su%3A%22An%C3%A1lisis%20vectorial%22}{Cálculo Vectorial de Colley}: sección 5.1, ejercicios 1--16 (impares); sección 5.2, ejercicios 1--22 (pares).

\paragraph{Pregunta de investigación:}
\begin{enumerate}[leftmargin=*]
    \item ¿Cómo se puede usar la integral doble para calcular la masa de una lámina plana con densidad variable?
    \item ¿Qué condiciones garantizan que se puede cambiar el orden de integración en una integral iterada?
\end{enumerate}

\paragraph{Para el próximo tema:}
Ver los videos:
\begin{itemize}[leftmargin=*]
    \item \href{https://youtu.be/wUF-lyyWpUc?si=bkW6BvU4wMgmxpjL}{Change of Variables \& The Jacobian}.
    \item \href{https://youtu.be/LUvynduoUX0?si=nv4TJNkUkNjgt_pV}{Change the order of integration to solve tricky integrals}.
    \item \href{https://youtu.be/IAb98ZgSJNw?si=RxvzpAw7ys28ySQw}{Intro to Polar Coordinates}.
\end{itemize}


\end{document}
